\section{Introduction}

Models for scientific images are routinely adapted across data collected by different experiments,
laboratories, scanners, or hospitals. These acquisition environments alter appearance while the
scientific prediction target remains nominally unchanged. A single dense feed-forward network must
represent every environment with the same parameters. Sparse mixture-of-experts (MoE) instead
supplies conditional capacity: each input activates only a subset of a larger parameter pool
\citep{fedus2022switch}. This
creates a plausible benefit---different adaptation demands need not interfere---and an equally
plausible failure---the router may partition data by acquisition nuisance and strand capacity when
a new environment appears.

Existing comparisons cannot answer which explanation is correct if an MoE is simply larger than
its dense baseline, if architecture and loss are tuned sequentially on the test domain, or if only
the winning configuration is analyzed. We therefore treat conditional capacity as the object of a
controlled kill test. We first require a pretrained transformer to demonstrate competent dense
full fine-tuning; a powerful model is not assumed to transfer merely because it is a foundation
model. We then change exactly one middle feed-forward block and hold the total parameter budget
fixed to the nearest realizable dense width. Only a seed-0 matched gain licenses replication or a
broader architectural study.

Our central question is: \emph{when adapting a pretrained vision model to scientific images
collected across acquisition environments, when does sparse conditional capacity improve
held-out-environment generalization over an equally large dense model?} We study RxRx1 cellular
microscopy, where 1,139 perturbations repeat across experimental batches and the official split
holds out complete experiments. The intended contributions are an auditable dense--sparse
comparison, a direct gradient-interference test, and mechanism analyses that distinguish reusable
conditional structure from acquisition nuisance.
